\subsection{Grounded-Laplacian route to an operator heavy-hitter locator}
\label{sec:grounded-laplacian-locator}

The remaining dynamic object is compatible with spectral vertex-sparsifier
machinery after a simple exact transformation.

\begin{proposition}[Grounded-Laplacian representation]
\label{prop:grounded-laplacian-representation}
Put $c_\alpha:=(1-\alpha)/2$ and
\begin{equation}
 H:=D^{1/2}QD^{1/2}=c_\alpha(D-A)+\alpha D.
 \label{eq:grounded-pagerank-matrix}
\end{equation}
Construct an augmented weighted graph on $V\cup\{\bot\}$ by giving every
original edge conductance $c_\alpha$ and adding a ground edge
$(v,\bot)$ of conductance $\alpha d_v$.  After pinning $\bot$, its Laplacian
principal minor is exactly $H$.

For every eliminated set $S$ and surviving set $W:=V\setminus S$,
\begin{equation}
 Q/S=D_W^{-1/2}(H/S)D_W^{-1/2}.
 \label{eq:schur-congruence-grounded}
\end{equation}
Consequently, for $A\subsetneq T$ and $v\notin T$,
\begin{equation}
 (Q/A)_{vv}-(Q/T)_{vv}
 =\frac{(H/A)_{vv}-(H/T)_{vv}}{d_v}.
 \label{eq:diagonal-loss-grounded}
\end{equation}
If a normalized-coordinate correction is written as
$\Delta_x=D_T^{1/2}\Delta_z$, then its normalized exterior response is
\begin{equation}
 \frac{(Q\Delta_x)_v}{\sqrt{d_v}}
 =\frac{(H\Delta_z)_v}{d_v},
 \qquad v\notin T.
 \label{eq:boundary-response-grounded}
\end{equation}
\end{proposition}

\begin{proof}
The first identity follows by direct multiplication.  The augmented graph
contributes $c_\alpha L$ on original edges and the diagonal
$\alpha D$ through its ground edges.  Diagonal congruence commutes with
principal Schur complementation: substituting
$Q_{UV}=D_U^{-1/2}H_{UV}D_V^{-1/2}$ into the block formula proves
\cref{eq:schur-congruence-grounded}.  Taking a diagonal difference proves
\cref{eq:diagonal-loss-grounded}.  Finally,
$Q\Delta_x=D^{-1/2}H\Delta_z$; divide its $v$th coordinate by
$\sqrt{d_v}$ to obtain \cref{eq:boundary-response-grounded}.
\end{proof}

\begin{proposition}[Exposed-incidence sufficiency]
\label{prop:exposed-incidence-sufficiency}
Fix a face $T$ and let
$W:=\{v\notin T:Q_{vT}\neq0\}$ be its graph boundary.  For every source set
$B\subseteq T$, all of the following objects are determined by
$Q_{TT}$, $Q_{WT}$, and $D_W$:
\begin{enumerate}[label=(\roman*)]
\item the degree-normalized response matrix
\begin{equation}
 \mathcal M_{W\leftarrow B}^{(T)}
 :=D_W^{-1/2}Q_{WT}Q_{TT}^{-1}E_B
   \left(E_B^TQ_{TT}^{-1}E_B\right)^{-1/2};
 \label{eq:exposed-normalized-response}
\end{equation}
\item every source-aware leverage and group mass formed from its row norms;
\item the normalized exterior demand caused by any dual source supported on
      $B$; and
\item every diagonal loss across a nested elimination $A\subsetneq T$, for
      labels in $W$.
\end{enumerate}
Consequently, changing arbitrary edges with both endpoints in $V\setminus T$
while keeping the degrees of labels in $W$ fixed changes none of these
objects.

In the adjacency-access model, scanning the rows of $T$ reveals every
internal edge, cut incidence, and boundary label.  One degree query for each
distinct label in $W$ completes the record.  Its size is
\begin{equation}
 O\!\left(|T|+\sum_{u\in T}d_u+|W|\right)
 =O\!\left(\operatorname{cvol}(T)\right).
 \label{eq:exposed-record-size}
\end{equation}
If a static operator-locator build on a record of size $M$ costs
$\widetilde O(MF)$ for a parameter factor $F$, rebuilding it only at anchor
faces whose charged volumes grow by a fixed constant factor costs
$\widetilde O(VF)$ through final exposed volume $V$.
\end{proposition}

\begin{proof}
The matrix in \cref{eq:exposed-normalized-response} is written entirely in
terms of the three declared blocks.  Its row norms give (ii), and multiplying
the same response by an arbitrary coefficient vector gives (iii).  For (iv),
\cref{eq:source-leverage-diagonal-loss} identifies every required diagonal
loss with the corresponding source-aware squared leverage.  Equivalently, it
is a diagonal of the response product in the block Schur formula.  No block
with both indices in $V\setminus T$ appears.

For the shared unweighted graph, $Q_{TT}$ is fixed by the degrees in $T$ and
the edges internal to $T$, while $Q_{WT}$ is fixed by the cut incidences and
the endpoint degrees.  Scanning the adjacency lists of $T$ reveals the first
two edge sets and all labels in $W$.  Moreover $|W|$ is at most the number of
cut incidences, hence at most $\sum_{u\in T}d_u$, proving
\cref{eq:exposed-record-size}.  The final claim is the geometric-series
argument of \cref{lem:geometric-rebuild} applied to the successive record
sizes.
\end{proof}

\begin{lemma}[Transposed fixed-anchor response sketch]
\label{lem:transposed-anchor-response-sketch}
Fix an anchor $A$, put $U:=V\setminus A$, and let
$\Pi\in\mathbb R^{p\times U}$ be any linear measurement matrix.  Define the
precomputed harmonic sketch
\begin{equation}
 Z_\Pi
 :=Q_{AA}^{-1}Q_{AU}D_U^{-1/2}\Pi^T
 \in\mathbb R^{A\times p}.
 \label{eq:transposed-harmonic-sketch}
\end{equation}
For any frontier $F\subseteq U$ and vector $y\in\mathbb R^F$, lift $y$ by
\begin{equation*}
 \Delta_A:=-Q_{AA}^{-1}Q_{AF}y,
 \qquad
 \Delta_F:=y,
 \qquad
 \Delta_{U\setminus F}:=0,
\end{equation*}
and put $q:=Q\Delta$.  Then $q|_A=0$ and
\begin{equation}
 \Pi D_U^{-1/2}q_U
 =\Pi D_U^{-1/2}Q_{UF}y-Z_\Pi^TQ_{AF}y.
 \label{eq:transposed-response-sketch}
\end{equation}
If $y$ is the exact Schur-frontier solution
\begin{equation*}
 \left(Q_{FF}-Q_{FA}Q_{AA}^{-1}Q_{AF}\right)y=h_F,
\end{equation*}
then $q_F=h_F$.  With $R:=U\setminus F$, the exterior-only measurement is
therefore
\begin{equation}
 \Pi_R D_R^{-1/2}q_R
 =\Pi D_U^{-1/2}q_U-\Pi_FD_F^{-1/2}h_F.
 \label{eq:frontier-contamination-subtraction}
\end{equation}

The two products $Q_{UF}y$ and $Q_{AF}y$ in
\cref{eq:transposed-response-sketch} are generated by scanning the rows of
$F$.  Thus, after the $p$ anchor right-hand sides in
\cref{eq:transposed-harmonic-sketch} have been solved and represented, the
identity requires no dense read of $\Delta_A$ and no pass over $U$.  It does
not declare the multiplication by the touched rows of $Z_\Pi$ free: their
storage, reads, sketch dimension $p$, and every anchor solve remain explicit
response charges.
\end{lemma}

\begin{proof}
The old-face block vanishes by construction.  On $U$,
\begin{equation*}
 q_U=Q_{UF}y-Q_{UA}Q_{AA}^{-1}Q_{AF}y.
\end{equation*}
Left-multiply by $\Pi D_U^{-1/2}$, use symmetry of $Q$, and transpose the
definition of $Z_\Pi$ to obtain
\cref{eq:transposed-response-sketch}.  On $F$, the same block multiplication
gives
\begin{equation*}
 q_F=\left(Q_{FF}-Q_{FA}Q_{AA}^{-1}Q_{AF}\right)y=h_F.
\end{equation*}
The measurement over $U$ is the sum of its $F$ and $R$ column restrictions,
so subtracting the known frontier contribution proves
\cref{eq:frontier-contamination-subtraction}.  Finally, every nonzero of
$Q_{VF}y$ is exposed by a row scan of $F$; the stated implementation ledger
follows directly from the two right-hand-side products.
\end{proof}

\begin{lemma}[Two-sided group sketch from transposed harmonic measurements]
\label{lem:two-sided-harmonic-group-sketch}
Use \cref{lem:transposed-anchor-response-sketch} with fixed anchor $A$,
frontier $F$, frontier Schur complement
\begin{equation*}
 K_F:=Q_{FF}-Q_{FA}Q_{AA}^{-1}Q_{AF},
\end{equation*}
and exterior $R:=V\setminus(A\cup F)$.  Let
$\mathcal U_F:\mathbb R^F\to\mathbb R^V$ be the corresponding frontier lift,
so that
\begin{equation*}
 (\mathcal U_Fy)_A=-Q_{AA}^{-1}Q_{AF}y,
 \qquad
 (\mathcal U_Fy)_F=y.
\end{equation*}
For a current hierarchy node $C\subseteq R$, define
\begin{equation}
 T_C:=D_C^{-1/2}Q_{CV}\mathcal U_FK_F^{-1/2},
 \qquad
 Z_F(C):=\|T_C\|_F^2.
 \label{eq:two-sided-harmonic-response}
\end{equation}
This is exactly the normalized source-aware group mass of $F$ relative to
$A$.

Draw independent standard Gaussian source vectors
$g_1,\ldots,g_r\in\mathbb R^F$ and, independently, global standard Gaussian
target vectors $a_1,\ldots,a_p\in\mathbb R^R$.  For every hierarchy node
$C$, let $a_{C,\ell}$ be the restriction of $a_\ell$ to $C$ and put
\begin{equation}
 \widetilde Z_F(C)
 :=\frac1{pr}\sum_{\ell=1}^p\sum_{t=1}^r
 \left(a_{C,\ell}^TT_Cg_t\right)^2.
 \label{eq:two-sided-harmonic-estimator}
\end{equation}
For $\xi,\delta\in(0,1)$ and a hierarchy with at most $M$ nodes, there is a
universal constant $c_{\rm 2sk}$ such that
\begin{equation}
 p,r\geq c_{\rm 2sk}\xi^{-2}\log(4M/\delta)
 \label{eq:two-sided-harmonic-count}
\end{equation}
implies, with probability at least $1-\delta$, simultaneously for every node
$C$,
\begin{equation}
 (1-\xi)^2Z_F(C)
 \leq\widetilde Z_F(C)
 \leq(1+\xi)^2Z_F(C).
 \label{eq:two-sided-harmonic-guarantee}
\end{equation}
Every scalar in \cref{eq:two-sided-harmonic-estimator} is an exterior
measurement in \cref{eq:frontier-contamination-subtraction}: choose one row of
$\Pi_R$ to be $a_{C,\ell}^TD_C^{-1/2}$, set
$y=K_F^{-1/2}g_t$, and subtract the known frontier contamination.  Therefore
the estimator requires no materialized exterior response and no scan over
$R$.
\end{lemma}

\begin{proof}
The source identity in
\cref{eq:normalized-frontier-source-identity} applied to $(A,F)$ shows that
$\mathcal U_FK_F^{-1/2}$ is the energy-normalized source-response basis, so
its normalized exterior row norms give the declared group mass.

Let $G:=[g_1\ \cdots\ g_r]$.  The Gaussian quadratic-form bound, unioned over
the hierarchy nodes, gives
\begin{equation*}
 (1-\xi)Z_F(C)
 \leq\frac1r\|T_CG\|_F^2
 \leq(1+\xi)Z_F(C)
\end{equation*}
simultaneously with failure probability at most $\delta/2$.  Condition on
$G$ and this event.  For $Y_C:=T_CG$, a second application of the same bound
to $Y_CY_C^T$, using the restricted target vectors $a_{C,\ell}$ and the remaining
failure probability $\delta/2$, gives
\begin{equation*}
 (1-\xi)\|Y_C\|_F^2
 \leq\frac1p\sum_{\ell=1}^p\|a_{C,\ell}^TY_C\|_2^2
 \leq(1+\xi)\|Y_C\|_F^2
\end{equation*}
for every node.  Combining the two displays proves
\cref{eq:two-sided-harmonic-guarantee}.  The final implementation statement
is exactly \cref{eq:frontier-contamination-subtraction} for the declared row
of $\Pi$ and source vector $y$.  The target restrictions are correlated
between nodes, but each restriction has the required standard Gaussian law;
the union bound uses no independence between nodes.  Sharing the global
targets therefore preserves the simultaneous guarantee.
\end{proof}

\begin{theorem}[Chebyshev whitening of a spectral frontier]
\label{thm:chebyshev-frontier-whitening}
Use the anchor, frontier, lift, and exact frontier Schur complement $K_F$ of
\cref{lem:two-sided-harmonic-group-sketch}.  Then
\begin{equation}
 \alpha I_F\preceq K_F\preceq I_F.
 \label{eq:frontier-schur-spectrum}
\end{equation}
Suppose an SPD frontier representation $\widetilde K_F$ satisfies
\begin{equation}
 (1-\upsilon)K_F\preceq\widetilde K_F
 \preceq(1+\upsilon)K_F,
 \qquad 0\leq\upsilon<1.
 \label{eq:frontier-spectral-approximation}
\end{equation}
Fix $\zeta\in(0,1/2)$.  There is a polynomial $p_q$ of degree
\begin{equation}
 q=O\!\left(
 \sqrt{\frac{1+\upsilon}{(1-\upsilon)\alpha}}\,
 \log\frac{1+\upsilon}{(1-\upsilon)\alpha\zeta}
 \right)
 =O_\upsilon\!\left(
 \frac1{\sqrt\alpha}\log\frac1{\alpha\zeta}
 \right)
 \label{eq:frontier-whitening-degree}
\end{equation}
such that, for
\begin{equation*}
 R_C:=D_C^{-1/2}Q_{CV}\mathcal U_F,
 \qquad
 Z_{F,q}(C):=\norm{R_Cp_q(\widetilde K_F)}_F^2,
\end{equation*}
every hierarchy node $C\subseteq R$ obeys
\begin{equation}
 a_-Z_F(C)\leq Z_{F,q}(C)\leq a_+Z_F(C),
 \qquad
 a_-:=\frac{(1-\zeta)^2}{1+\upsilon},
 \quad
 a_+:=\frac{(1+\zeta)^2}{1-\upsilon}.
 \label{eq:frontier-whitening-sandwich}
\end{equation}

Replace $K_F^{-1/2}g_t$ in
\cref{eq:two-sided-harmonic-estimator} by
$y_t:=p_q(\widetilde K_F)g_t$ and denote the resulting estimator by
\begin{equation}
 \widetilde Z_{F,q}(C)
 :=\frac1{pr}\sum_{\ell=1}^p\sum_{t=1}^r
 \left(a_{C,\ell}^TR_Cy_t\right)^2.
 \label{eq:chebyshev-whitened-group-estimator}
\end{equation}
Under \cref{eq:two-sided-harmonic-count}, with probability at least
$1-\delta$, simultaneously for every hierarchy node,
\begin{equation}
 (1-\xi)^2a_-Z_F(C)
 \leq\widetilde Z_{F,q}(C)
 \leq(1+\xi)^2a_+Z_F(C).
 \label{eq:chebyshev-whitened-group-guarantee}
\end{equation}
Consequently
\begin{equation}
 \widehat Z_{F,q}(C)
 :=\frac{\widetilde Z_{F,q}(C)}{(1-\xi)^2a_-}
 \label{eq:chebyshev-whitened-one-sided}
\end{equation}
is a simultaneous one-sided oracle with
\begin{equation}
 Z_F(C)\leq\widehat Z_{F,q}(C)
 \leq\Gamma_{\xi,\zeta,\upsilon}Z_F(C),
 \qquad
 \Gamma_{\xi,\zeta,\upsilon}
 :=
 \left(\frac{1+\xi}{1-\xi}\right)^2
 \left(\frac{1+\zeta}{1-\zeta}\right)^2
 \frac{1+\upsilon}{1-\upsilon}.
 \label{eq:chebyshev-whitened-oracle-factor}
\end{equation}

If one multiplication by $\widetilde K_F$ costs $T_{\rm mv}$ work, all
$r$ source vectors $y_t$ are prepared once, before any adaptive group search,
in
\begin{equation}
 O(rqT_{\rm mv})
 =\widetilde O_{\xi,\zeta,\upsilon,\delta}\!\left(
  \frac{T_{\rm mv}}{\sqrt\alpha}
 \right)
 \label{eq:frontier-whitening-work}
\end{equation}
work.  In the exact case $\widetilde K_F=K_F$, one frontier multiplication
\begin{equation}
 K_Fy=Q_{FF}y-Q_{FA}Q_{AA}^{-1}Q_{AF}y
 \label{eq:exact-frontier-matvec}
\end{equation}
uses one charged anchor-response application and
$O(\operatorname{cvol}(F))$ additional incidence work.  Thus neither an exact
frontier square root nor an exact solve for every Gaussian source is required.
\end{theorem}

\begin{proof}
For $y\in\mathbb R^F$, the Schur variational identity gives
\begin{equation*}
 y^TK_Fy
 =\min_{x\in\mathbb R^A}
 \begin{pmatrix}x\\y\end{pmatrix}^{\!T}
 Q_{A\cup F,A\cup F}
 \begin{pmatrix}x\\y\end{pmatrix}
 \geq\alpha\norm{y}_2^2.
\end{equation*}
Also $K_F\preceq Q_{FF}\preceq I$, proving
\cref{eq:frontier-schur-spectrum}.  Hence the spectrum of
$\widetilde K_F$ lies in
$[(1-\upsilon)\alpha,1+\upsilon]$.  Chebyshev approximation on this interval
supplies the degree in \cref{eq:frontier-whitening-degree} and the relative
inverse-square-root bound
\begin{equation*}
 (1-\zeta)^2\widetilde K_F^{-1}
 \preceq p_q(\widetilde K_F)^2
 \preceq(1+\zeta)^2\widetilde K_F^{-1}.
\end{equation*}
Inverting \cref{eq:frontier-spectral-approximation} gives
\begin{equation*}
 \frac1{1+\upsilon}K_F^{-1}
 \preceq\widetilde K_F^{-1}
 \preceq\frac1{1-\upsilon}K_F^{-1}.
\end{equation*}
Congruence by $R_C$ and taking traces proves
\cref{eq:frontier-whitening-sandwich}.

Apply the two Gaussian concentration steps in the proof of
\cref{lem:two-sided-harmonic-group-sketch} to the fixed matrix
$R_Cp_q(\widetilde K_F)$ in place of $T_C$.  They bound
\cref{eq:chebyshev-whitened-group-estimator} between
$(1-\xi)^2Z_{F,q}(C)$ and $(1+\xi)^2Z_{F,q}(C)$ simultaneously for all
nodes.  Combining this with \cref{eq:frontier-whitening-sandwich} proves
\cref{eq:chebyshev-whitened-group-guarantee,eq:chebyshev-whitened-oracle-factor}.

A Chebyshev recurrence uses $q$ multiplications by $\widetilde K_F$ for each
source vector, proving \cref{eq:frontier-whitening-work}.  Finally,
\cref{eq:exact-frontier-matvec} follows from the definition of $K_F$.
Both sparse products incident on $F$ are formed by scanning its rows; the
middle operation is exactly one application of the stored anchor response.
\end{proof}

\begin{theorem}[Coded bucket locator for row energy]
\label{thm:coded-bucket-row-locator}
Let $\mathcal W$ be a set of labels with an injective binary encoding in a
universe of size $N\geq2$, and let
$X\in\mathbb R^{\mathcal W\times r}$.  Fix $\theta>0$, $K\geq1$, and
$\delta\in(0,1)$ such that
\begin{equation}
 \|X\|_F^2\leq K\theta.
 \label{eq:coded-locator-mass-cap}
\end{equation}
Put
\begin{equation*}
 L:=\left\lceil\log_2N\right\rceil,
 \qquad
 B:=\left\lceil128KL\right\rceil,
 \qquad
 R:=\left\lceil\log_4(N/\delta)\right\rceil.
\end{equation*}
For every repetition $s\in[R]$, independently draw a pairwise-independent
uniform hash $h_s:\mathcal W\to[B]$ and pairwise-independent Rademacher
signs $\sigma_s(v)$.  For $b\in[B]$ form the vector measurements
\begin{align}
 M_{s,b,0}
 &:=\sum_{v:h_s(v)=b}\sigma_s(v)X_{v,:},
 \label{eq:coded-locator-total-bucket}\\
 M_{s,b,\ell}
 &:=\sum_{\substack{v:h_s(v)=b\\
              \operatorname{bit}_\ell(v)=1}}
       \sigma_s(v)X_{v,:},
 \qquad 1\leq\ell\leq L.
 \label{eq:coded-locator-bit-bucket}
\end{align}
Decode one label from every pair $(s,b)$ by setting its $\ell$th bit to one
when
\begin{equation}
 \|M_{s,b,\ell}\|_2>
 \|M_{s,b,0}-M_{s,b,\ell}\|_2,
 \label{eq:coded-locator-bit-test}
\end{equation}
breaking ties arbitrarily, and retain the decoded label only when it belongs
to $\mathcal W$.  Let $\mathcal C$ be the union of the retained labels.  Then,
with probability at least $1-\delta$,
\begin{equation}
 \left\{v\in\mathcal W:\|X_{v,:}\|_2^2\geq\theta\right\}
 \subseteq\mathcal C,
 \qquad
 |\mathcal C|\leq BR.
 \label{eq:coded-locator-output}
\end{equation}
The complete measurement family has
\begin{equation}
 p=B(L+1)R
 =O\!\left(K\log^2N\log(N/\delta)\right)
 \label{eq:coded-locator-row-count}
\end{equation}
scalar target rows.  It is implicit: a label contributes to exactly one total
bucket and at most $L$ bit buckets in each repetition.  The decoder uses
$O(pr)$ arithmetic and produces at most
$O(K\log N\log(N/\delta))$ candidates before exact validation.
\end{theorem}

\begin{proof}
Fix a heavy label $v$, one repetition $s$, and its bucket
$b:=h_s(v)$.  For each bit $\ell$ and value $c\in\{0,1\}$, let
\begin{equation*}
 N_{\ell,c}
 :=\sum_{\substack{u\neq v:\ h_s(u)=b\\
              \operatorname{bit}_\ell(u)=c}}
       \sigma_s(u)X_{u,:}.
\end{equation*}
Pairwise independence of the signs removes all cross terms, and pairwise
independence of the hash gives
\begin{equation*}
 \mathbb E\|N_{\ell,c}\|_2^2
 \leq\frac{\|X\|_F^2}{B}
 \leq\frac{K\theta}{B}.
\end{equation*}
Since $\|X_{v,:}\|_2^2\geq\theta$, Markov's inequality yields
\begin{equation*}
 \Pr\!\left\{
  \|N_{\ell,c}\|_2>\tfrac14\|X_{v,:}\|_2
 \right\}
 \leq\frac{16K}{B}.
\end{equation*}
A union bound over the $2L$ pairs $(\ell,c)$ shows that all of these noise
vectors have norm at most $\|X_{v,:}\|_2/4$ with probability at least
$1-32KL/B\geq3/4$.

On that event, if the $\ell$th bit of $v$ is one, then
$M_{s,b,\ell}$ contains $\sigma_s(v)X_{v,:}$ plus noise of norm at most one
quarter of the signal, whereas
$M_{s,b,0}-M_{s,b,\ell}$ is noise of norm at most one quarter of the signal.
Their norms are therefore at least three quarters and at most one quarter of
$\|X_{v,:}\|_2$, respectively.  The inequalities reverse when the bit is
zero.  Thus every bit test in
\cref{eq:coded-locator-bit-test} is correct and the bucket decodes $v$.

The repetitions are independent, so the probability that none decodes this
fixed $v$ is at most $4^{-R}\leq\delta/N$.  A union bound over at most $N$
heavy labels proves the containment in
\cref{eq:coded-locator-output}.  There is at most one retained candidate per
pair $(s,b)$, proving the size bound.  Counting the total and bit rows proves
\cref{eq:coded-locator-row-count}.  Finally, evaluating all vector norms in
the bit tests costs $O(pr)$, and the stated column sparsity follows directly
from \cref{eq:coded-locator-total-bucket,eq:coded-locator-bit-bucket}.
\end{proof}

\begin{corollary}[One fixed-event shared coded harmonic bank for a packed query]
\label{cor:shared-coded-harmonic-bank}
Use the anchor $A$, frontier $F$, and full exterior graph boundary
\begin{equation*}
 W:=\{v\notin A\cup F:Q_{v,A\cup F}\neq0\}.
\end{equation*}
Let $\widetilde K_F$ be the spectral frontier and use the Chebyshev source vectors
$y_t=p_q(\widetilde K_F)g_t$ of
\cref{thm:chebyshev-frontier-whitening}.  Assume $F\neq\varnothing$, and let
the exposed labels have identifiers in a universe of size $N\geq2$.  For
$r$ independent source probes define the row matrix
\begin{equation}
 X_{v,t}:=\frac{(Q\mathcal U_Fy_t)_v}{\sqrt{r d_v}},
 \qquad v\in W,\qquad 1\leq t\leq r.
 \label{eq:coded-harmonic-row-matrix}
\end{equation}
Fix $\xi,\delta\in(0,1/2)$ and take
\begin{equation*}
 r\geq c_{\rm sk}\xi^{-2}\log(4M/\delta),
\end{equation*}
where $M$ bounds the current hierarchy nodes and leaves.  For a correction
energy radius $E>0$ and normalized margin $m>0$, put
\begin{equation}
 \tau:=\frac{m^2}{E^2},
 \qquad
 \theta:=(1-\xi)a_-\tau,
 \qquad
 K:=\max\!\left\{1,
  \frac{(1+\xi)a_+|F|}{(1-\xi)a_-\tau}
 \right\},
 \label{eq:coded-harmonic-budget}
\end{equation}
with $a_-,a_+$ from
\cref{eq:frontier-whitening-sandwich}.  Draw the coded bank of
\cref{thm:coded-bucket-row-locator}, with failure parameter $\delta/2$,
independently after the current
$A,F,\widetilde K_F$ are fixed.  With probability at least $1-\delta$, its
candidate set contains every potentially ambiguous leaf
\begin{equation}
 \left\{v\in W:Z_F(\{v\})\geq\frac{m^2}{E^2}\right\},
 \label{eq:coded-harmonic-ambiguous-set}
\end{equation}
and has size
\begin{equation}
 \widetilde O_{\xi,\zeta,\upsilon,\delta}\!\left(
  1+|F|\frac{E^2}{m^2}
 \right).
 \label{eq:coded-harmonic-candidate-count}
\end{equation}
The number of target rows in the single shared bank obeys the same bound up
to logarithms.  Exact validation of those candidates therefore returns the
same ambiguous set as the hierarchy search, without constructing one target
harmonic row per reached hierarchy node.

More precisely, define every signed bucket/bit coefficient on the whole
identifier universe $U:=V\setminus A$ and use the resulting rows as $\Pi$ in
\cref{lem:transposed-anchor-response-sketch}; the locator uses their
restriction $\Pi_W$.  All entries of $\Pi_WX$ are supplied exactly by the
transposed harmonic identity after subtracting the known frontier block.
Labels outside $F\cup W$ contribute zero because they are not adjacent to the
support of the lifted vector.  Thus the harmonic bank may be formed at the
fixed anchor before $W$ is known.  Its capacity may depend on the already-fixed
public event parameters $F,E,m,\xi,\zeta,\upsilon$ and the assigned failure
budget, but it may not be selected after inspecting sampled source responses or
bank answers.  A block built earlier must use a declared capacity envelope and
an assignment rule independent of its unobserved contents.  Let
$\mathcal C_{\rm harm}(A,\Pi,Y)$ denote every charged anchor right-hand-side
construction, anchor solve or update, representation write, touched old-face
row read, and harmonic application needed for the $p$ target rows and the
source matrix $Y=[y_1\ \cdots\ y_r]$.  Apart from this declared response
charge, the query costs
\begin{equation}
 \widetilde O\!\left(
  \frac{T_{\rm mv}}{\sqrt\alpha}
  +rRL\operatorname{cvol}(F)
  +pr
 \right)
 +\mathcal C_{\rm val}(\mathcal C),
 \label{eq:coded-harmonic-query-work}
\end{equation}
where $R,L,p$ are from
\cref{thm:coded-bucket-row-locator} and
$\mathcal C_{\rm val}$ charges construction or maintenance of the exposed
boundary-membership dictionary, membership lookup, exact leaf evaluation, and
output.  In addition to retained response state, the explicit measurement,
decoder, and candidate arrays use $O(pr+BR)$ transient cells.  The
$rRL\operatorname{cvol}(F)$ term forms the direct bucket
contributions by incidence scans: each touched label updates only one bucket
and its $L$ bit counters per repetition.
Forming the anchor right-hand sides does not enumerate $U$: only columns with
$Q_{Av}\neq0$ contribute, and their identifiers and degrees are already in
the exposed anchor-cut record of
\cref{prop:exposed-incidence-sufficiency}.

For an adaptive trace, the same conclusion is valid event by event only if,
after the event tuple $(A,F,\widetilde K_F,E,m)$ is fixed, the event either
draws both fresh Gaussian source probes $g_t$ (hence $Y$) and a fresh
hash/sign target bank, or consumes an unused, unobserved block containing both
random layers and independent of the complete history.  All such blocks are
charged in $\mathcal C_{\rm harm}$.  A simultaneous whole-trace guarantee
additionally assigns summable budgets $\sum_j\delta_j\leq\delta$ and uses
\begin{equation*}
 r_j\geq c_{\rm sk}\xi^{-2}\log(4M_j/\delta_j),
 \qquad
 R_j=\left\lceil\log_4(2N/\delta_j)\right\rceil,
\end{equation*}
with every resulting logarithm included in the aggregate block charge.
Reusing either random layer after its answers have influenced a later event is
not justified by this corollary.
\end{corollary}

\begin{proof}
Condition on the current frontier.  The first Gaussian quadratic-form step in
\cref{lem:two-sided-harmonic-group-sketch}, unioned over the current leaves
and the root group $W$, gives, with failure probability at most $\delta/2$,
\begin{equation*}
 (1-\xi)a_-Z_F(\{v\})
 \leq\|X_{v,:}\|_2^2
 \leq(1+\xi)a_+Z_F(\{v\})
\end{equation*}
for every $v\in W$, and
\begin{equation*}
 \|X\|_F^2
 \leq(1+\xi)a_+Z_F(W)
 \leq(1+\xi)a_+|F|.
\end{equation*}
The last inequality is the source-aware leverage packing bound
\cref{eq:source-aware-normalized-mass-total}.  Hence
\cref{eq:coded-locator-mass-cap} holds with the $\theta,K$ in
\cref{eq:coded-harmonic-budget}, and every leaf in
\cref{eq:coded-harmonic-ambiguous-set} has sampled row energy at least
$\theta$.

Condition on this source event and draw the independent coded bank with
failure budget $\delta/2$.  Apply
\cref{thm:coded-bucket-row-locator}.  A final union bound proves the
no-false-negative claim.  Substitution of
\cref{eq:coded-harmonic-budget} into
\cref{eq:coded-locator-output,eq:coded-locator-row-count} proves the candidate
and target-row bounds.

For a source column $y_t$, the corresponding coded measurements are exactly
\begin{equation*}
 \frac1{\sqrt r}\Pi_WD_W^{-1/2}
      (Q\mathcal U_Fy_t)_W.
\end{equation*}
This is \cref{eq:frontier-contamination-subtraction}: the full measurement is
formed first and its known $F$ contribution is subtracted.  The
source-preparation cost is
\cref{eq:frontier-whitening-work}.  A scan of the frontier rows forms the
direct exterior terms; the implicit bank has at most $R(L+1)$ nonzeros per
label, so all $r$ source columns cost the second term of
\cref{eq:coded-harmonic-query-work}.  The decoder and its vector norms cost
$O(pr)$.  Every operation involving the anchor representation or final
candidate is, by definition, retained in the two explicit charges.  Finally,
  both random layers in a fresh unused block are independent of the history
  that fixes the current event, so the conditional union-bound argument of
\cref{thm:causal-orthogonal-response-sketches} applies event by event.
\end{proof}

The corollary closes the measurement-count question for one packed query:
the target rows are output-sensitive even at full cut rank and are shared by
all would-be hierarchy nodes.  It does not close the graph-local theorem,
because $\mathcal C_{\rm harm}$ may still include too many anchor solves,
stored old-face reads, or independent bank blocks.  The next lemma must bound
that aggregate charge through an epoch, or replace the harmonic bank by the
positive shifted-debt realization below.

\begin{theorem}[Certified adaptive reuse of multiscale coded banks]
\label{thm:certified-adaptive-coded-bank-reuse}
Fix an identifier universe of size $N\geq2$ and an upper bound $J_{\max}$ on
the number of packed queries in one fixed-anchor epoch.  At exact reference
query $j$, let
\begin{equation*}
 X_j\in\mathbb R^{\mathcal W_j\times r_j},
 \qquad
 \|X_j\|_F^2\leq K_j\theta_j,
 \qquad
 K_j\geq1,\quad \theta_j>0,
\end{equation*}
where $\mathcal W_j$ is any current subset of the identifier universe.  The
row-heavy superset is
\begin{equation}
 \mathcal G_j
 :=\{v\in\mathcal W_j:\|(X_j)_{v,:}\|_2^2\geq\theta_j\}
 \label{eq:certified-reference-row-heavy-set}
\end{equation}
and suppose the exact reference event set $\mathcal H_j$ satisfies
$\mathcal H_j\subseteq\mathcal G_j$.  Exact validation returns
$\mathcal H_j$ from any candidate set containing it, and the reference
controller's next state is a fixed function of this validated set and of
randomness independent of the target banks.  Rejected false candidates,
their order, and the raw bucket values do not change this state.

Put $\kappa_\ell:=2^\ell$.  For every capacity level that is used, draw one
target bank $\Pi^{(\ell)}$ from
\cref{thm:coded-bucket-row-locator} with mass parameter
$K=\kappa_\ell$ and failure parameter $\delta/J_{\max}$.  At query $j$, use
the smallest level $\ell(j)$ with $K_j\leq\kappa_{\ell(j)}$, decode its
candidate set, exactly validate every candidate, and update the controller
only with the validated set.  The same bank at a level is reused at every
later query assigned to that level.

Then, with probability at least $1-\delta$, no reference-heavy label is
missed at any query and the adaptive run follows exactly the reference state
trace.  If
\begin{equation*}
 K_{\max}:=\max_{j\leq J_{\max}}K_j,
\end{equation*}
the total number of distinct target rows over all capacity levels is
\begin{equation}
 P_{\rm epoch}
 =O\!\left(
  K_{\max}\log^2N\log(NJ_{\max}/\delta)
 \right),
 \label{eq:certified-multiscale-bank-rows}
\end{equation}
not the sum of the per-query row counts.  Query $j$ decodes and validates at
most
\begin{equation}
 O\!\left(
  K_j\log N\log(NJ_{\max}/\delta)
 \right)
 \label{eq:certified-multiscale-query-candidates}
\end{equation}
candidates, up to harmless rounding at $K_j=1$.

The levels need not be allocated in advance.  A level may be drawn lazily
the first time its capacity is requested, after the current query tuple is
fixed and before its measurements are inspected.  This requires no advice
about $K_{\max}$ and preserves the same row bound.

For the harmonic realization in
\cref{cor:shared-coded-harmonic-bank}, this proves that one harmonic target
bank per realized capacity level suffices throughout a certified adaptive
epoch.  The source probes may remain fresh with summable failure budgets as in
\cref{thm:causal-orthogonal-response-sketches}; on their simultaneous success
event, the truly ambiguous leaves are the required sets $\mathcal H_j$ and
are contained in the sampled row-heavy sets $\mathcal G_j$.  Their failure
probability is added separately.  The target-side response charge has the
exact decomposition
\begin{equation}
 \mathcal C_{\rm harm}^{\rm epoch}
 =\mathcal C_{\rm build}
   \bigl(A,\{\Pi^{(\ell)}\}_{\ell}\bigr)
  +\sum_j\mathcal C_{\rm apply}
   \bigl(A,\Pi^{(\ell(j))},Y_j\bigr),
 \label{eq:certified-harmonic-build-apply-split}
\end{equation}
where the first term builds each distinct row only once and the second term
charges every later old-face read and application.  Thus adaptive
correctness no longer forces one fresh harmonic bank build per event.  The
theorem does not bound either implementation term in
\cref{eq:certified-harmonic-build-apply-split}.
\end{theorem}

\begin{proof}
First condition on the graph, the exact validation rule, and all randomness
other than the target banks.  This fixes the exact reference trace
recursively, including every
$(\mathcal W_j,X_j,\theta_j,K_j,\mathcal G_j,\mathcal H_j)$.  For a fixed reference query
$j$, its assigned bank has capacity
$\kappa_{\ell(j)}\geq K_j$, so
\begin{equation*}
 \|X_j\|_F^2\leq\kappa_{\ell(j)}\theta_j.
\end{equation*}
By \cref{thm:coded-bucket-row-locator}, the probability that this bank misses
a member of $\mathcal G_j$, and hence the probability that it misses a member
of $\mathcal H_j\subseteq\mathcal G_j$, is at most
$\delta/J_{\max}$.  The events may be
correlated because a level is reused; a union bound over the fixed reference
queries still gives total failure probability at most $\delta$.

It remains to justify using the reference queries in an adaptive execution.
If the execution ever leaves the reference trace, let $j$ be the first such
query.  Every earlier candidate set contained its complete reference event
set, and exact validation deleted every false candidate.  By the controller
hypothesis, all earlier state transitions therefore equal their reference
transitions.  Hence the query tuple at $j$ is exactly the fixed reference
tuple, and divergence can occur only if its assigned bank misses a member of
$\mathcal H_j$.  The probability of any first divergence is consequently
bounded by the preceding union.  This is the certified first-failure
coupling; no for-all adaptive-sketch guarantee is assumed.

For one level, substitute the common failure parameter
$\delta/J_{\max}$ into
\cref{eq:coded-locator-row-count}.  Since the realized capacities are powers
of two,
\begin{equation*}
 \sum_{\ell:\ \kappa_\ell\leq2K_{\max}}\kappa_\ell
 <4K_{\max}.
\end{equation*}
This proves \cref{eq:certified-multiscale-bank-rows}.  Minimality of
$\ell(j)$ gives $\kappa_{\ell(j)}<2K_j$ unless $K_j=1$; apply
\cref{eq:coded-locator-output} to prove
\cref{eq:certified-multiscale-query-candidates}.

For lazy allocation, condition on the certified history through the first
request to a new level.  The current query tuple is already fixed, while the
new bank is fresh and independent.  The same first-failure argument applies
from that point onward.  Geometric summation is unchanged.  Finally, the
global identifier construction in
\cref{cor:shared-coded-harmonic-bank} lets the same $\Pi^{(\ell)}$ be
restricted to every later boundary.  Building each harmonic extension once
and charging every application separately gives
\cref{eq:certified-harmonic-build-apply-split}.
\end{proof}

The theorem removes the adaptive fresh-bank multiplier from
$\mathcal C_{\rm harm}$.  The remaining high-rank question is now strictly
implementation-side: represent the $P_{\rm epoch}$ harmonic rows and apply
the selected level after each light event within the desired aggregate old-face
ledger.

\begin{theorem}[Bidirectional harmonic application]
\label{thm:bidirectional-harmonic-application}
Fix an anchor $A$, put $U:=V\setminus A$, and let
$\Pi\in\mathbb R^{p\times U}$ be any target bank.  For a frontier
$F\subseteq U$ and a source matrix $Y\in\mathbb R^{F\times r}$, put
\begin{equation*}
 W:=Q_{AF}Y,
 \qquad
 \mathsf H(\Pi,Y)
 :=\Pi D_U^{-1/2}Q_{UA}Q_{AA}^{-1}W.
\end{equation*}
The old-face term in every column of
\cref{eq:transposed-response-sketch} has the two exactly equal realizations
\begin{equation}
 \mathsf H(\Pi,Y)
 =\left(Q_{AA}^{-1}Q_{AU}D_U^{-1/2}\Pi^T\right)^TW
 =\Pi D_U^{-1/2}Q_{UA}X,
 \qquad Q_{AA}X=W.
 \label{eq:bidirectional-harmonic-identity}
\end{equation}
Thus the target orientation solves $p$ anchor right-hand sides once and
stores their harmonic extensions, whereas the source orientation solves the
$r$ current right-hand sides and stores no target harmonic extension.

Suppose now that $\Pi$ is one signed hash-and-bit bank from
\cref{thm:coded-bucket-row-locator}, and let $\sigma$ be the number of its
nonzeros in one identifier column.  The source orientation can form
$\mathsf H(\Pi,Y)$ exactly, after solving for $X$, by one scan of the
anchor--exterior cut in
\begin{equation}
 O\!\left(pr+\sigma r\,\operatorname{cvol}(A)\right)
 \label{eq:source-oriented-coded-bank-work}
\end{equation}
arithmetic and hash operations.  Including the direct term and the known
frontier subtraction in
\cref{eq:transposed-response-sketch,eq:frontier-contamination-subtraction}
adds
$O(\sigma r\,\operatorname{cvol}(F))$ work.  The transient arrays use
$O(|A|r+pr)$ cells in addition to the exact multi-right-hand-side solve
state.  All degree reads are those of the exposed-incidence record in
\cref{prop:exposed-incidence-sufficiency}.  In particular, this realization
has no harmonic-bank build or hidden scan of the exterior, but it pays one
old-face cut scan and one anchor solve for every source block.
\end{theorem}

\begin{proof}
Symmetry of $Q_{AA}$ gives
\begin{align*}
 \left(Q_{AA}^{-1}Q_{AU}D_U^{-1/2}\Pi^T\right)^TW
 &=\Pi D_U^{-1/2}Q_{UA}Q_{AA}^{-1}W\\
 &=\Pi D_U^{-1/2}Q_{UA}X,
\end{align*}
which proves \cref{eq:bidirectional-harmonic-identity}.

For the work claim, initialize the $p\times r$ accumulator to zero.  While
scanning a cut incidence $(a,v)$ with $a\in A$ and $v\in U$, its contribution
to source column $t$ is
\begin{equation*}
 \frac{Q_{va}X_{a,t}}{\sqrt{d_v}}\,\Pi_{:,v}.
\end{equation*}
The identifier $v$ determines the $\sigma$ nonzero coordinates and signs of
$\Pi_{:,v}$, so the incidence costs $O(\sigma r)$.  There are at most
$\sum_{a\in A}d_a$ such incidences.  Initializing the accumulator costs
$O(pr)$, proving \cref{eq:source-oriented-coded-bank-work}.  Scanning $F$
and applying the same column rule forms the direct exterior term and the
frontier subtraction at the stated additional charge.  No other exterior
coordinate is read.
\end{proof}
